% sec:schw-kerr-schild — Kerr-Schild Form
\section{Kerr-Schild Form}
\label{sec:schw-kerr-schild}

The Eddington--Finkelstein coordinates of the previous section remove the
coordinate singularity at $r = 2M$, but they introduce a new problem: the
metric is written in spherical coordinates $(v, r, \theta, \varphi)$, which
have their own coordinate singularity at the poles ($\theta = 0, \pi$), and
the flat-spacetime limit yields $\eta_{\mu\nu}$ in spherical rather than
Cartesian form.  The ray tracer needs Cartesian spatial coordinates to set
up initial conditions and to avoid special treatment at the poles.  The
Kerr--Schild form delivers both: a single Cartesian chart that is
horizon-penetrating, pole-free, and whose inverse metric costs nothing more
than a sign flip.

\paragraph{From Eddington--Finkelstein to Kerr--Schild.}

The ingoing Eddington--Finkelstein metric (\cref{eq:ef-ingoing}) uses the
advanced time $v = t + r_*$ as its time coordinate.  Define a new time
coordinate by subtracting the radial coordinate:
%
\begin{equation}
  \bar{t} = v - r \,.
  \label{eq:ks-time-def}
\end{equation}
%
Since $v = t + r_*$ and $r_* = r + 2M\ln(r/2M - 1)$
(\cref{eq:tortoise,eq:ef-v-def}), the Kerr--Schild time differs from
Schwarzschild time outside the horizon ($r > 2M$) by a logarithmic
correction: $\bar{t} = t + 2M\ln(r/2M - 1)$.  The defining relation
$\bar{t} = v - r$ remains valid at and inside the horizon, where
Schwarzschild time is undefined.

Substituting $\d v = \d\bar{t} + \d r$ into the Eddington--Finkelstein line
element \eqref{eq:ef-ingoing} and writing $f = 2M/r$,
%
\begin{align}
  \d s^2 &= -(1 - f)\,(\d\bar{t} + \d r)^2
    + 2\,(\d\bar{t} + \d r)\,\d r
    + r^2\,\d\Omega^2 \notag \\
  &= \underbrace{-\d\bar{t}^2 + \d r^2
    + r^2\,\d\Omega^2}_{\text{Minkowski}}
    \;+\; f\,(\d\bar{t} + \d r)^2 \,.
  \label{eq:ks-spherical}
\end{align}
%
The first line substitutes directly; the second separates the
flat-spacetime contribution from a gravitational correction.  The entire
effect of curvature appears in the second term: a distortion along the
radial null direction $\d\bar{t} + \d r$, with strength $f = 2M/r$ that
falls off as the Newtonian potential.  (Expanding
$-\d\bar{t}^2 + \d r^2 + f(\d\bar{t} + \d r)^2$ reproduces the
coefficients $(-1+f)$, $2f$, $(1+f)$ for
$\d\bar{t}^2$, $\d\bar{t}\,\d r$, $\d r^2$, confirming the
factorisation.)  The one-form $\d\bar{t} + \d r$ has components
$l_\mu = (1, 1, 0, 0)$ in coordinates $(\bar{t}, r, \theta, \varphi)$,
so the metric takes the Kerr--Schild form
%
\begin{equation}
  \metric = \eta_{\mu\nu} + f\,l_\mu\,l_\nu \,,
  \qquad f = \frac{2M}{r} \,,
  \label{eq:ks-schw-metric}
\end{equation}
%
fulfilling the decomposition previewed in \cref{sec:gr-exact-solutions}:
the full curved metric is a flat background plus a rank-one correction
built from a single covector $l_\mu$ and a scalar function $f$.
\histnote{Kerr and Schild introduced this decomposition in 1965
\cite{KerrSchild1965}, originally for the rotating Kerr metric.
The Schwarzschild case is a special limit that was recognised almost
immediately.}

Converting to Cartesian spatial coordinates
$(x, y, z) = (r\sin\theta\cos\varphi,\, r\sin\theta\sin\varphi,\,
r\cos\theta)$ replaces $\d r^2 + r^2\,\d\Omega^2$ with
$\d x^2 + \d y^2 + \d z^2$, and the radial one-form becomes
$\d r = (x\,\d x + y\,\d y + z\,\d z)/r$.  The components of $l_\mu$ in
Cartesian Kerr--Schild coordinates $(\bar{t}, x, y, z)$ are
%
\begin{equation}
  l_\mu = \left(1,\; \frac{x}{r},\; \frac{y}{r},\;
    \frac{z}{r}\right) , \qquad r = \sqrt{x^2 + y^2 + z^2} \,.
  \label{eq:ks-null-oneform}
\end{equation}
%
These are the Kerr--Schild data stated in \cref{sec:gr-exact-solutions},
now derived from first principles.

\paragraph{The null property and the inverse metric.}

The covector $l_\mu$ is null with respect to the flat metric:
%
\begin{equation}
  \eta^{\mu\nu}\,l_\mu\,l_\nu = -(1)^2 + \frac{x^2 + y^2 + z^2}{r^2}
  = -1 + 1 = 0 \,.
  \label{eq:ks-flat-null}
\end{equation}
%
The payoff of this null property is immediate and substantial.  For a
generic perturbation $\metric = \eta_{\mu\nu} + h_{\mu\nu}$, the inverse
metric requires a Neumann series $\inversemetric = \eta^{\mu\nu} -
h^{\mu\nu} + h^{\mu\alpha}h_\alpha{}^\nu - \cdots$ that may not even
converge near the horizon.  For a null rank-one perturbation, the series
terminates after one term.  Writing
$l^\mu = \eta^{\mu\alpha}l_\alpha = (-1,\, x/r,\, y/r,\, z/r)$ and
computing
%
\begin{equation}
  g^{\mu\alpha}\,g_{\alpha\nu}
  = \delta^\mu{}_\nu
    + f\,l^\mu\,l_\nu
    - f\,l^\mu\,l_\nu
    - f^2\,l^\mu\underbrace{(l^\alpha\,l_\alpha)}_{=\;0}\,l_\nu
  = \delta^\mu{}_\nu \,,
  \label{eq:ks-inverse-verify}
\end{equation}
%
confirms that the inverse metric is
%
\begin{equation}
  \inversemetric = \eta^{\mu\nu} - f\,l^\mu\,l^\nu \,.
  \label{eq:ks-inverse}
\end{equation}
%
The two cross terms in \cref{eq:ks-inverse-verify} are each
$f\,l^\mu\,l_\nu$ and cancel directly, while the $f^2$ term vanishes
because $l_\mu$ is null.  The inverse differs from the forward metric by
nothing more than a sign flip in front of the correction.
\warnnote{The sign flip between $\metric = \eta_{\mu\nu} +
f\,l_\mu\,l_\nu$ and $\inversemetric = \eta^{\mu\nu} -
f\,l^\mu\,l^\nu$ relies entirely on $l_\mu$ being null.  For a
non-null perturbation, the inverse involves the full Neumann series;
the null property truncates it after one term.}

\paragraph{Explicit components.}

Writing $n^i = x^i/r$ for the Euclidean unit radial vector ($i = 1, 2, 3$;
in Cartesian coordinates $n_i = n^i$ since spatial indices are raised and
lowered with $\delta_{ij}$), the Kerr--Schild metric components in
coordinates $(\bar{t}, x, y, z)$ are
%
\begin{equation}
  g_{\bar{t}\bar{t}} = -1 + f \,, \quad
  g_{\bar{t}i} = f\,n_i \,, \quad
  g_{ij} = \delta_{ij} + f\,n_i n_j \,.
  \label{eq:ks-components}
\end{equation}
%
The pattern carries over to the inverse: the correction term flips sign
while the mixed time--space components are unchanged.
%
\begin{equation}
  g^{\bar{t}\bar{t}} = -1 - f \,, \quad
  g^{\bar{t}i} = f\,n^i \,, \quad
  g^{ij} = \delta^{ij} - f\,n^i n^j \,.
  \label{eq:ks-inv-components}
\end{equation}
%
At $r = 2M$ the function $f = 1$: every component is finite, and the
metric is perfectly well-behaved.  At the ISCO ($r = 6M$), $f = 1/3$ ---
the metric deviates from flat spacetime by roughly $30\%$, yet every
component remains a simple rational function of the coordinates.  At large
$r$, $f \to 0$ and the metric smoothly approaches $\eta_{\mu\nu}$ in
Cartesian form --- precisely the flat-spacetime coordinates the codebase
uses for initial conditions.

The metric determinant is worth noting for numerical work: $\det(g) = -1$,
independent of position.  This follows from the matrix determinant lemma:
$\det(A + uv^T) = (1 + v^T\!A^{-1}u)\det(A)$.  With $A = \eta$,
$u = f\,l$, and $v = l$, the correction factor is
$1 + f\,\eta^{\mu\nu}l_\mu\,l_\nu = 1$ since $l_\mu$ is null.
\physnote{The unit determinant means the natural volume element
$\detmetric\;\d^4\!x$ reduces to the coordinate volume
$\d\bar{t}\,\d x\,\d y\,\d z$ --- volume integrals over the Kerr--Schild
chart have the same form as in flat spacetime.}

\paragraph{Computational advantages.}

Three features make the Kerr--Schild form the natural choice for the
geodesic integrator:
%
\begin{enumerate}
\item \textbf{Horizon-penetrating:} All metric components are finite
  at $r = 2M$, so the integrator steps smoothly through the horizon
  without adaptive refinement or coordinate switching.
\item \textbf{Pole-free:} Cartesian coordinates have no
  coordinate singularity at $\theta = 0, \pi$, eliminating special
  treatment of polar geodesics.
\item \textbf{Algebraically simple inverse:} The sign-flip formula
  \eqref{eq:ks-inverse} means the inverse metric is as cheap to
  evaluate as the forward metric --- no matrix inversion at each
  integration step.
\end{enumerate}
%
\coderef{Spacetime}{SchwarzschildKS}
The codebase implements these properties directly.  The type
\texttt{SchwarzschildKS} stores only the mass $M$; the polymorphic function
\texttt{schwarzschildMetricFn} computes all ten independent components
of $\metric$ from \cref{eq:ks-components} for any \texttt{Floating}
type, enabling automatic differentiation for Christoffel symbols without
hand-coded derivatives (\cref{ch:autodiff}).
\implnote{The polymorphic metric function accepts \texttt{Double} for
production ray tracing and AD types for exact derivative computation.
A single code path serves both purposes.}
%
\begin{lstlisting}[language=Haskell]
schwarzschildMetricFn :: Floating a => a -> [a] -> [a]
schwarzschildMetricFn m [_, cx, cy, cz] =
  let r  = sqrt (cx*cx + cy*cy + cz*cz)
      h  = 2*m / r
      lx = cx / r; ly = cy / r; lz = cz / r
  in [ -1 + h,   h*lx,       h*ly,       h*lz
     ,           1 + h*lx*lx, h*lx*ly,    h*lx*lz
     ,                        1 + h*ly*ly, h*ly*lz
     ,                                     1 + h*lz*lz ]
\end{lstlisting}
%
The ten entries are the independent components of the symmetric
$4 \times 4$ metric, listed in upper-triangular order.  Each line
of the code maps directly to \cref{eq:ks-components}: the \texttt{-1 + h}
entry is $g_{\bar{t}\bar{t}} = -1 + f$, the off-diagonal entries
\texttt{h*lx}, \texttt{h*ly}, \texttt{h*lz} are $g_{\bar{t}i} = f\,n_i$,
and the spatial block is $\delta_{ij} + f\,n_i n_j$.

The Kerr--Schild decomposition is not specific to the Schwarzschild
spacetime.  The Kerr metric for a rotating black hole admits the same
$\metric = \eta_{\mu\nu} + f\,l_\mu\,l_\nu$ structure, with $f$ and
$l_\mu$ generalised to incorporate the spin parameter $a$; setting $a = 0$
recovers \cref{eq:ks-schw-metric} exactly.  This universality allows the
codebase to use the same integration infrastructure for both spacetimes.
\xrefnote{The Kerr--Schild form of the Kerr metric is derived in
\cref{sec:kerr-ks}, where $l_\mu$ acquires a twist from the
black hole's rotation.}
